% GENERATED, edit the French unit or the translation table
\mxnewpage\section*{EV3-A -- Mock exam 2}\label{EV3-A}
\addcontentsline{toc}{section}{EV3-A -- Mock exam 2}
\textbf{Paper reserved for the teacher}. Programme studied in chapters 1 to 10.\par Duration: 2 h. No calculator. Marked out of 20. 1--3 April 2027.
\subsection*{Part 1 -- Fluency (6 points)}
Twelve questions, \num{0,5} point each. Only one correct answer; no justification required. An incorrect or missing answer scores 0, with no loss of points.
\begin{enumerate}
\item Calculate $3+2\times3$.\par A. $9$\quad B. $15$\quad C. $8$\quad D. $12$
\item Calculate $\frac12+\frac13$.\par A. $\frac{2}{5}$\quad B. $\frac{5}{6}$\quad C. $\frac{1}{6}$\quad D. $\frac{5}{3}$
\item Calculate $\dfrac{2^3}{2^{1}}$.\par A. $4$\quad B. $1$\quad C. $8$\quad D. $2$
\item A value of 50 increases by \pc{20}. Final value?\par A. $10$\quad B. $40$\quad C. $60$\quad D. $70$
\item A value of 80 is obtained after a \pc{20} decrease. Initial value?\par A. $96$\quad B. $60$\quad C. $100$\quad D. $64$
\item A \pc{25} increase then a \pc{20} decrease: overall coefficient?\par A. $\frac{4}{5}$\quad B. $\frac{21}{20}$\quad C. $\frac{9}{10}$\quad D. $1$
\item The solution of $2x+3=0$ is\par A. $\frac{3}{2}$\quad B. $\frac{-3}{2}$\quad C. $-6$\quad D. $6$
\item The slope between $\pt{0;3}$ and $\pt{2;9}$ is\par A. $6$\quad B. $-3$\quad C. $2$\quad D. $3$
\item The mean of $3, 5, 7$ is\par A. $6$\quad B. $15$\quad C. $3$\quad D. $5$
\item A sector of $90^\circ$ represents what fraction of the total?\par A. $\frac{1}{3}$\quad B. $\frac{1}{4}$\quad C. $\frac{3}{4}$\quad D. $\frac{1}{2}$
\item A probability is $\dfrac{2}{5}$. Its complement is\par A. $\frac{3}{5}$\quad B. $\frac{5}{2}$\quad C. $\frac{2}{5}$\quad D. $\frac{1}{5}$
\item For $f(x)=x^2+3$, $f(-2)$ is\par A. $12$\quad B. $7$\quad C. $5$\quad D. $-1$
\end{enumerate}
\mxnewpage\subsection*{Part 2 -- Independent exercises (14 points)}
Every step must be justified. Exact results are expected unless stated otherwise.
\mxneed{12}
\subsection*{Exercise 1 -- Lottery and dispersion (5 points)}
A lottery pays a gross gain $G$ of 0, 4 or 12 monetary units with respective probabilities $\dfrac12$, $\dfrac38$ and $\dfrac18$. To take part, one pays 3 units. Let $X=G-3$ be the net gain.
\begin{enumerate}
\item Give the distribution of the net gain X. \hfill(1 pt)
\item Calculate the expectation and variance of X. \hfill(2 pts)
\item The gross gains and the entry fee are doubled. What become the mean and standard deviation of the net gain? \hfill(2 pts)
\end{enumerate}
\mxneed{12}
\subsection*{Exercise 2 -- Derivative and horizontal tangent (5 points)}
The function $f$ is defined on $\R$ by $f(x)=(2x-3)e^x$. We study its extremum and the tangent to its curve at the point with x-coordinate 0. All values are expected in exact form.
\begin{enumerate}
\item Calculate $f(0)$, then $f\prime(x)$ and $f\prime(0)$. \hfill(1 pt)
\item Study the variation of $f$ on $\R$ and give its exact extremum. \hfill(2 pts)
\item Determine the equation of the tangent at the point with x-coordinate 0. State whether its point of contact is above or below the x-axis. \hfill(2 pts)
\end{enumerate}
\mxneed{12}
\subsection*{Exercise 3 -- Angle and length (4 points)}
In an orthonormal frame, given the points $A\pt{0;0}$, $B\pt{4;0}$ and $C\pt{2;2\sqrt3}$. We look for the nature of the triangle $ABC$ by comparing lengths and scalar product.
\begin{enumerate}
\item Calculate AB and AC. \hfill(1 pt)
\item Determine the cosine then the angle BAC in radians. \hfill(1 pt)
\item Determine BC and the nature of the triangle. \hfill(2 pts)
\end{enumerate}
